\vspace{-1mm}\section{Related Work}\vspace{-1mm}
Chain-of-Thought (CoT) elicits intermediate rationales, improving multi-step reasoning on math and logic tasks. Scaling short CoT to Long CoT is difficult: coherence often degrades, and cold-start gains frequently require targeted training or high-quality trajectories~\citep{guo2025deepseek,moshkov2025aimo}. 
A common approach distills stepwise solutions from strong teacher models into weaker students~\citep{muennighoff2025s1,zhao20251}. 
Outcomes hinge on reasoning quality: strong reasoning models transfer useful behaviors~\citep{ye2025limo}, whereas weaker instruction models may mimic format without robust Long CoT capability~\citep{du2025the,chandra2025shape}.

Long CoT training elicits three behaviors: deep reasoning, self-reflection, and self-exploration~\citep{o1,guo2025deepseek,chen2025towards}. Work has studied their roles; \citet{madaan2023self} used self-reflection to revise earlier steps, and \citet{shinn2023reflexion} combined reflection and exploration to improve robustness. Early studies framed CoT as sequences of these behaviors, emphasizing step-level imitation and local coherence~\citep{wei2022chain,kojima2022large,chen2024unlocking}. Later work used tree- or graph-structured reasoning to capture branching and revisitation~\citep{yao2023tree,yao2024got,hu2024treeplanner,besta2024graph}.

Although trees or graphs represent individual Long CoT traces by modeling behaviors as nodes, they do not capture the overall distribution of logical behaviors. In contrast, our approach models Long CoT as a molecular-like structure, with edges encoding stable distributions of reasoning behaviors, to test how their arrangement and interactions support effective learning.

